\documentclass[sigconf,nonacm]{acmart}

\settopmatter{printacmref=false}

\usepackage{amsmath}
% amssymb omitted: acmart's math setup already defines AMS symbols (\Bbbk etc.); loading amssymb again triggers ``already defined'' errors.
\usepackage{booktabs}
\usepackage{graphicx}
\usepackage{url}
\usepackage{microtype}
\usepackage{array}
\usepackage{tcolorbox}
\tcbuselibrary{listings,breakable,skins}
% Avoid writing colored prompt titles into PDF bookmarks (math, \\, etc.).
\tcbset{bookmark=false}
\definecolor{DefnBoxGray}{gray}{0.94}
% Courier (pcr): avoids inconsolata (zi4) missing italic in listings.
\newcommand{\PromptMono}{\footnotesize\fontencoding{T1}\fontfamily{pcr}\selectfont}
\lstdefinestyle{promptlist}{
  basicstyle=\PromptMono,
  columns=fullflexible,
  breaklines=true,
  breakautoindent=false,
  tabsize=2,
  showstringspaces=false,
  frame=none,
  aboveskip=0pt,
  belowskip=0pt,
  lineskip=0pt,
}
\lstdefinestyle{promptlistdark}{
  basicstyle=\color{white}\PromptMono,
  columns=fullflexible,
  breaklines=true,
  breakautoindent=false,
  tabsize=2,
  showstringspaces=false,
  frame=none,
  aboveskip=0pt,
  belowskip=0pt,
  lineskip=0pt,
}
\newcommand{\AppendixSnippet}[1]{appendix_snippets/#1}
\newenvironment{PromptSystemBox}[1]{%
  \begin{tcolorbox}[
    width=\linewidth,
    enhanced, arc=2mm, outer arc=2mm,
    colback=black, colframe=black, colupper=white,
    coltitle=white, colbacktitle=black,
    fonttitle=\bfseries\sffamily\footnotesize,
    title={#1}, titlerule=0pt, boxrule=0pt,
    top=1mm, boxsep=3mm,
    before upper={\lstset{linewidth=\linewidth}},
    before skip=\baselineskip, after skip=\baselineskip,
    breakable]
}{\end{tcolorbox}}
\newenvironment{PromptUserBox}[1]{%
  \begin{tcolorbox}[
    width=\linewidth,
    enhanced, arc=2mm, outer arc=2mm,
    colback=white, colframe=black, colupper=black,
    coltitle=black, colbacktitle=white,
    fonttitle=\bfseries\sffamily\small,
    title={#1}, titlerule=0.55pt, titlerule style=black,
    boxrule=0.85pt, top=0.5mm, boxsep=3.5mm, left=3.5mm, right=3.5mm,
    before upper={\lstset{linewidth=\linewidth}},
    before skip=\baselineskip, after skip=\baselineskip,
    breakable]
}{\end{tcolorbox}}
\newenvironment{DefnBox}[1]{%
  \begin{tcolorbox}[
    width=\linewidth,
    enhanced, arc=2mm,
    colback=DefnBoxGray, colframe=black!55,
    fonttitle=\bfseries\footnotesize,
    title={#1},
    boxrule=0.55pt, boxsep=3mm, breakable]
}{\end{tcolorbox}}
\newcommand{\LstPrompt}[1]{\lstinputlisting[style=promptlist]{\AppendixSnippet{#1}}}
\newcommand{\LstPromptDark}[1]{\lstinputlisting[style=promptlistdark]{\AppendixSnippet{#1}}}

\graphicspath{{figures/}{results/figures/}}

\newcommand{\missingfigure}[2]{%
  \fbox{%
    \begin{minipage}[c][0.25\textheight][c]{0.92\linewidth}
      \centering
      Figure asset not found: \texttt{\detokenize{#1}}\\[0.5em]
      #2
    \end{minipage}%
  }%
}
\newcommand{\maybeincludegraphics}[2]{%
  \IfFileExists{#1}{%
    \Description{#2}%
    \includegraphics[width=\linewidth]{#1}%
  }{%
    \IfFileExists{results/#1}{%
      \Description{#2}%
      \includegraphics[width=\linewidth]{results/#1}%
    }{%
      \missingfigure{#1}{#2}%
    }%
  }%
}

\title[Conflict as Structural Stress in LLM Cross-Personality Rewriting]{Conflict as Structural Stress in LLM Cross-Personality Rewriting}

\author{Kaixun Wang}
\author{Xuyang Liu}
\affiliation{
  \institution{Southern University of Science and Technology}
  \city{Shenzhen}
  \country{China}
}

\begin{document}

\begin{abstract}
We frame cross-persona conflict as \textbf{structural stress}, not as a uniform ``hybrid vigor'' boost: outcomes \textbf{split by dimension}---automatic novelty (combined) rises with conflict curvature ($\beta_{d^2}{>}0$, Table~\ref{tab:quadratic}), while value and coherence bend downward ($\beta_{d^2}{<}0$), and headline creativity sits between those forces ($\beta_{d^2}{=}{-}0.111$, $p{\approx}4.6{\times}10^{-7}$, $n{=}4620$).
The pre-registered T3-only fit places a quadratic vertex at $d^\ast{\approx}0.47$ on discrete-main T3 rows (Table~\ref{tab:h1}), signalling \textbf{medium} conflict within that subset; pooled binned means of $C_{\mathrm{auto}}$ vs.\ $d_H$ nonetheless trend broadly downward at scale (Figure~\ref{fig:main}), so we do \emph{not} describe headline creativity as a classic rise-then-fall ``inverted-U'' in the aggregate plot.
The pooled sample comprises \textbf{4716} rows (two main arms, mechanism \textbf{60}, multihop 96). Exploratory bins tie rising conflict to higher mean sentiment displacement and non-monotone sentence-order similarity (Figure~\ref{fig:variability}). On the mechanism arm, mean $C_{\mathrm{auto}}$ ranks \textbf{M1}${>}$\textbf{M3}${>}$\textbf{M2}${>}$\textbf{M0} on fifteen high-$d$ cells per mode (descriptive).
Separately, open-source \textbf{Qwen3-4B/8B/14B} replications on the same discrete T3 protocol each show significantly negative quadratic curvature on $C_{\mathrm{auto}}$ vs.\ $d_H$ under the \textbf{same Gaussian mixed-effects specification} as Table~\ref{tab:scale_regression} (genre fixed effects with source-level random intercepts; $n{=}720$ T3 rows per checkpoint: $\beta_{d^2}{=}{-}0.641$, $-0.744$, $-0.432$, respectively, all $p{<}10^{-3}$).
Figure~\ref{fig:scale_inverted_u} is a \textbf{visual supplement only}: it overlays genre-averaged \emph{OLS} predicted curves on the same Qwen rows to sketch shape and hybrid bands; those OLS vertices are not re-used as the primary inferential estimand and need not coincide numerically with the MixedLM $d^\ast$ column in Table~\ref{tab:scale_regression}.
The judge-composite analogue appears in Figure~\ref{fig:scale_inverted_u_judge}.
A data-driven \textbf{Space-L} (SBERT$\to$PCA, Procrustes to Space-H) recovers the same $S{\times}R$ layout (Figure~\ref{fig:space_l_validation}); H7 regressions on $d_L$ keep negative $\beta_{d^2}$ at all Qwen checkpoints, so main analyses use $d_H$.
\end{abstract}

\keywords{AI writing, persona prompting, creativity evaluation, style transfer, structural stress, LLM-as-judge, mixed-effects regression}

\maketitle

\section{Scope of This Report}

Experiments span corpus-scale rewriting with automatic metrics and LLM-as-judge scores, merged for pooled analysis.
The pooled proprietary sample comprises \textbf{4716} rows across discrete main, continuous main, mechanism, and multihop arms (Table~\ref{tab:coverage}).
The open-source Qwen3 replication contributes an additional $3{\times}1140=3420$ rows (one metrics file per checkpoint) analysed in the scale-specific tables and figures.
Mixed-effects regressions use all rows with finite conflict scalar $d_H$ and finite outcomes on each dependent variable (hence $n{=}4620$ for headline automatic creativity on the proprietary pooled regression sample, fewer than 4716 when some rows lack finite $d_H$ or outcomes).

\paragraph{Interpretive frame (structural stress).}
Rather than treating conflict as a homogeneous creativity elevator, we emphasize \textbf{dimensional splitting}: the same stressor can elevate novelty-like signals while eroding value and coherence support for $C_{\mathrm{auto}}{=}N_{\mathrm{auto}}{\cdot}V_{\mathrm{auto}}$, consistent with creativity theories that distinguish generative fluency from evaluative constraint~\cite{boden} and with dual-process accounts that separate rapid associative production from slower monitoring of consistency and judgment~\cite{kahneman}.
Under this reading, an interior optimum near $d^\ast{\approx}0.47$ (Table~\ref{tab:h1}) marks a \textbf{medium-conflict} regime where rewriting remains productive before coherence/value penalties steepen---not an endorsement of maximal persona distance.

\section{Experimental Design}

\noindent\textbf{Appendices.}
Full rewrite and judge prompt templates, automatic-metric formulas, supplementary regression/binning robustness tables, and the exploratory structural-collapse GEE are in Appendices~\ref{app:prompts}--\ref{app:collapse}.

\subsection{Persona Space and Conflict}

Space-H maps each source to $(S,R)\in[-1,1]^2$: increasing $S$ toward $+1$ indicates more analytic, System-2-like prose (toward $-1$, more intuitive and affective); increasing $R$ toward $+1$ indicates a more adventurous register, toward $-1$ a more conservative one. Persona prototypes carry fixed coordinates, and rewrite rows inherit conflict from the source--target pair.

\paragraph{Mapping sources into Space-H.}
The default pipeline uses \textbf{deterministic HuggingFace probes} (no LLM sampling): $S$ clips $(f{-}e)$ where $f$ is a RoBERTa formality probability and $e$ is non-neutral emotion mass from a DistilRoBERTa emotion classifier; $R$ is the dot product of a StyleDistance sentence embedding with a fixed unit axis loaded from precomputed axis vectors (so register contrast is geometric rather than hand-labelled). An alternate configuration uses an LLM coordinate scorer that returns JSON $(S,R)$ on a 1--5 scale, averaged over repeated seeds with spread retained as a stability diagnostic---both backends normalize to $[-1,1]$ before conflict computation.

\paragraph{Reliability and genre structure.}
Coordinates are screened with a \textbf{genre-separation diagnostic}: per-genre means and standard deviations of $S$ and $R$, pairwise Euclidean distances between genre centroids in $(S,R)$, and one-way ANOVA plus Kruskal--Wallis tests for differences in $S$ and $R$ across genres. These checks justify treating Space-H as an ordinal scaffold for conflict rather than as a literal psychometric scale; absolute distances should still be interpreted cautiously.
Figure~\ref{fig:coord_space_h} plots all corpus sources in $(S,R)$ by genre under the same deterministic mapping used for conflict construction.

\begin{figure}[t]
\centering
\maybeincludegraphics{figures/coord_space_h_distribution.png}{Missing figure.}
\caption{Source-level Space-H coordinates $(S,R)$ by genre ($n{=}60$ sources). Axis directions follow the project persona specification: $S$ runs from System-1 (intuitive, affective) at $-1$ to System-2 (analytic, rational) at $+1$; $R$ runs from conservative register at $-1$ to adventurous register at $+1$.}
\label{fig:coord_space_h}
\end{figure}

Scalar conflict $d_H$ and directional $\Delta_H=(\Delta S,\Delta R)$ are computed from these vectors (Euclidean $d_H$ after the project's normalization).

\subsection{Space-L validation (learned embedding space)}

\textbf{Space-L} tests whether the hand-crafted Space-H corners are recoverable from text: per persona we generate $\geq 200$ reference passages (GPT-4o / GPT-4o mini), embed with SBERT, PCA to $k{=}2$, and Procrustes-align to the four H persona vectors (\texttt{build\_space\_l}).
Figure~\ref{fig:space_l_validation} summarises the result.
\textbf{Left:} aligned persona layouts in H vs.\ L (Procrustes disparity $0.017$; cross-generator disparity $0.004$).
\textbf{Right:} H7 robustness---T3 MixedLM $\beta_{d^2}$ on $C_{\mathrm{auto}}$ stays negative for $d_H$ (Table~\ref{tab:scale_regression}) and for learned conflict $d_L$ on all Qwen3 checkpoints ($n{=}720$ each).
PC1 tracks register $R$ ($|r|{=}0.995$), PC2 tracks $S$ ($r{=}0.989$); PCA explains $\approx$88\% centroid variance.
Main regressions, prompts, and CrEval matching therefore use \textbf{$d_H$}; Space-L is a confirmatory geometry check, not a second primary estimand.

\begin{figure}[t]
\centering
\maybeincludegraphics{figures/space_l_validation.png}{Space-L Procrustes alignment and H7 curvature comparison.}
\caption{Space-L validation. \emph{Left:} four persona prototypes after Procrustes alignment (RC/EC/RA/EA). \emph{Right:} T3-only MixedLM $\beta_{d^2}$ on $C_{\mathrm{auto}}$ using $d_H$ vs.\ $d_L$ (Qwen3; all $p{<}10^{-3}$ on $d_L$). Built from cached \texttt{.cache/space\_l/}; reproduce with \texttt{scripts/plot\_space\_l\_validation.py}.}
\label{fig:space_l_validation}
\end{figure}

\subsection{Main Arms: Discrete and Continuous Targets}

The discrete main arm uses T0 (identity), T1 (nearest persona), T2 (random persona placebo), and T3 (directed cross-personality).
A parallel continuous-target arm repeats the same structure with continuously sampled T3 targets; both arms enter pooled regressions and genre summaries, while condition means in Table~\ref{tab:condition} use the discrete arm only.
Figure~\ref{fig:d_distribution_discrete_continuous} contrasts the induced T3 distributions of scalar conflict $d_H$: discrete targets concentrate on a small set of corner-persona distances, whereas continuous targets fill the unit disk and yield lower typical $d_H$ (the same comparison yields a handful of descriptive summary statistics in the supplementary materials).

\begin{table}[t]
\centering
\small
\begin{tabular}{lrrrr}
\toprule
Arm & Rows & T0 & T1/T2 & T3 or modes \\
\midrule
Discrete main & 2280 & 120 & 360 / 360 & 1440 T3 \\
Continuous main & 2280 & 120 & 360 / 360 & 1440 T3 \\
Mechanism & 60 & -- & -- & 15 each M0--M3 \\
Multihop & 96 & -- & -- & 96 hop rows \\
\midrule
Total pooled & 4716 & 240 & 1440 / 1440 & 2880 T3, 60 mech, 96 mh.\\
\bottomrule
\end{tabular}
\caption{Row counts by experimental arm. The Qwen3 replication arm contributes a separate $3{\times}1140=3420$ rows analysed in model-specific scale tables only, not in these pooled-arm totals.}
\label{tab:coverage}
\end{table}

\begin{figure}[t]
\centering
\maybeincludegraphics{figures/d_distribution_discrete_vs_continuous.png}{Missing figure.}
\caption{T3 conflict distance $d_H$ under discrete persona corners versus continuously sampled targets in Space-H ($n{=}1440$ T3 rows per main arm). Overlaid density histograms with a shared bin count.}
\label{fig:d_distribution_discrete_continuous}
\end{figure}

\paragraph{Open-source small-model replication arm.}
Three \textbf{Qwen3} checkpoints (4B, 8B, 14B) are served through a local deployment and run the \emph{same} discrete main-arm protocol (T0--T3, joint mode, identical prompts and row schema) as the proprietary-generator arms; all downstream metric computation, judging, merging, and analysis use the same pipeline so that cross-scale comparisons isolate generator capacity rather than preprocessing differences. Each Qwen checkpoint contributes 1140 rows (T0=60, T1=180, T2=180, T3=720), with the T3 block driving the regressions in Tables~\ref{tab:scale_regression}--\ref{tab:judge_auto_divergence}.

\section{Measurement Model}

\subsection{Automatic novelty $N_{\mathrm{auto}}$}

$N_{\mathrm{auto}}$ is the equal-weight mean of three \textbf{genre-relative} channels, each mapped to $[0,1]$ after comparing the rewrite to statistics of \emph{source texts in the same genre}: (i)~distinct-2 type-token ratio vs.\ genre mean/variance; (ii)~bigram novelty rate vs.\ typical intra-genre overlap; (iii)~optional SBERT-based $1{-}$cosine distance vs.\ mean pairwise embedding distance within the genre (this channel is informative only when embedding computation is enabled). Z-score--style clipping stabilizes outliers before rebinding to $[0,1]$. The design deliberately removes genre-wide verbosity confounds that would otherwise make poetry appear ``novel'' by construction.

\subsection{Automatic value $V_{\mathrm{auto}}$}

Let entailment be the probability of natural-language inference \emph{entailment} from source premise to generated hypothesis (cross-encoder MNLI). Let coherence be GPT-2 perplexity of the rewrite mapped monotonically to $[0,1]$ (inverse perplexity proxy for fluency). The primary value scalar is the \textbf{arithmetic mean} of entailment and coherence. A geometric mean variant $\sqrt{e\cdot c}$ is computed for robustness tables only.

\subsection{Automatic creativity $C_{\mathrm{auto}}$}

Primary automatic creativity is the \textbf{product}
\[
  C_{\mathrm{auto}} = N_{\mathrm{auto}}\, V_{\mathrm{auto}},
\]
clipped so both factors lie in $[0,1]$ beforehand. Multiplication encodes a conjunctive requirement: high scores need both divergence from the source \emph{and} retention of propositional/fluency support---consistent with definitions that treat creativity as combining originality with acceptable usefulness~\cite{runco}. An additive composite would reward dominance on one axis alone and blur interpretation when novelty and value move against each other under conflict.

Pooling therefore lets novelty and value/coherence move in opposite directions under rising $d_H$, producing the split documented in Table~\ref{tab:quadratic}.

\subsection{LLM-as-judge (validity arm)}

Absolute judging uses a fixed rubric: seven Likert dimensions (novelty, surprise, imagery, fidelity, fluency, consistency, logical completeness) with strict JSON outputs on a 1--5 scale; each generation is scored independently relative to the source and genre. \textbf{Dual judges} are used in this study---OpenAI GPT-4o and a non-reasoning DeepSeek-Chat endpoint (DeepSeek thinking mode explicitly disabled at the API layer to keep all tokens in the JSON message and to eliminate the empty-content failures that were observed on the earlier V4-Pro reasoning route)---spanning distinct provider families rather than repeated calls to the same endpoint. Per-dimension scores are averaged across judges with successful parses. A pairwise creative-comparison prompt is implemented for head-to-head evaluation with randomized option order, but the merged analysis emphasises the absolute branch for stable row-wise fields. Aggregated judge novelty pools novelty-type dimensions; weak correlation with $C_{\mathrm{auto}}$ at the row level ($r{\approx}0.20$; Table~\ref{tab:validity_corr}) signals partial construct divergence between human-like rubric novelty and the automatic NLI$\times$perplexity definition rather than ``failure'' of either side in isolation.

Judge fields are merged into the analysis table for validity checks; they \textbf{do not} replace $C_{\mathrm{auto}}$ as the prespecified primary outcome.

\subsection{Judge validity and construct checks}

Table~\ref{tab:validity_corr} reports Pearson correlations on rows with merged dual-judge fields.
On the proprietary discrete-main file ($n{=}2280$), the headline composites align moderately: $r(C_{\mathrm{auto}},C_{\mathrm{judge}}){=}0.43$; rubric novelty tracks automatic novelty at $r{=}0.46$; $C_{\mathrm{auto}}$ correlates with automatic novelty and value at $r{\approx}0.66$ each; $r(C_{\mathrm{auto}},\mathrm{novelty}_{\mathrm{judge}}){=}0.20$.
On the pooled Qwen3 replication ($n{=}3420$), $r(C_{\mathrm{auto}},C_{\mathrm{judge}}){=}0.31$ with the same weak-to-moderate $r(C_{\mathrm{auto}},\mathrm{novelty}_{\mathrm{judge}}){=}0.20$.
Thus the two \emph{composite} creativity channels move together positively at the row level, even though their \textbf{quadratic} responses to $d_H$ can point in opposite directions (Table~\ref{tab:judge_auto_divergence}).
Figure~\ref{fig:aux} (left) visualises sub-score correlations on the experiment-wide flat table (judge sub-dimensions vs.\ automatic components; the heatmap does not include $C_{\mathrm{judge}}$ itself).
Judge fidelity vs.\ automatic novelty is negative on proprietary rows ($r{\approx}{-}0.41$), consistent with a novelty--fidelity tension; judge coherence and fidelity remain positively associated ($r{\approx}0.60$).

\begin{table}[t]
\centering
\small
\begin{tabular}{lrrrr}
\toprule
Sample & $n$ & $r(C_{\mathrm{auto}},C_{\mathrm{judge}})$ & $r(C_{\mathrm{auto}},\mathrm{nov.}_{\mathrm{judge}})$ & $r(N_{\mathrm{auto}},\mathrm{nov.}_{\mathrm{judge}})$ \\
\midrule
Proprietary discrete main & 2280 & 0.43 & 0.20 & 0.46 \\
Qwen3 pooled (4B/8B/14B) & 3420 & 0.31 & 0.20 & 0.41 \\
Qwen3-4B / 8B / 14B & 1140 ea. & 0.21 / 0.35 / 0.31 & 0.19 / 0.22 / 0.18 & 0.32 / 0.43 / 0.46 \\
\bottomrule
\end{tabular}
\caption{Pearson validity correlations on merged judge rows (all conditions in each metrics file). T3-only $r(C_{\mathrm{auto}},C_{\mathrm{judge}})$ is 0.36 (proprietary) and 0.29 (Qwen3 pooled).}
\label{tab:validity_corr}
\end{table}

To stress-test whether this channel dependence extends beyond linear correlations, we repeated the \textbf{T3-only} scale-style regressions used alongside Table~\ref{tab:scale_regression}, replacing $C_{\mathrm{auto}}$ with the judge-based composite $C_{\mathrm{judge}}$ (computed from merged judge scores in the same analysis pipeline).
Each fit uses $C_{\cdot} \sim d_H + d_H^2 + C(\mathrm{genre})$ with a source-level random intercept on T3 rows (Qwen checkpoints: $n{=}720$ each; proprietary OpenAI runs pooled across two model assignments: $n{=}1440$ T3 discrete rows in one fit, as in Table~\ref{tab:judge_auto_divergence}).

\begin{table}[t]
\centering
\small
\begin{tabular}{lrrrrrr}
\toprule
Generator & $\beta_{d^2}^{\mathrm{auto}}$ & $p(d^2)$ & Sig. & $\beta_{d^2}^{\mathrm{judge}}$ & $p(d^2)$ & Sig. \\
\midrule
Qwen3-4B  & $-0.641$ & $1.7{\times}10^{-9}$ & Yes & $-0.049$ & $0.851$ & No \\
Qwen3-8B  & $-0.744$ & $3.0{\times}10^{-8}$ & Yes & $+0.226$ & $0.429$ & No \\
Qwen3-14B & $-0.432$ & $1.2{\times}10^{-4}$ & Yes & $+0.700$ & $5.0{\times}10^{-3}$ & Yes \\
GPT-4o    & $-0.618$ & $7.0{\times}10^{-10}$ & Yes & $+1.007$ & $1.8{\times}10^{-5}$ & Yes \\
\bottomrule
\end{tabular}
\caption{T3-only quadratic curvature for automatic vs.\ judge creativity by generator (Gaussian mixed-effects with genre fixed effects and source random intercepts). All four generators show significant negative curvature on $C_{\mathrm{auto}}$. On $C_{\mathrm{judge}}$, Qwen3-14B and the proprietary GPT pool show significantly positive curvature, opposite the automatic channel.}
\label{tab:judge_auto_divergence}
\end{table}

On the Qwen3 arms---where the same downstream judging stack is shared across checkpoints---$C_{\mathrm{judge}}$ need not track $C_{\mathrm{auto}}$.
For \textbf{Qwen3-4B}, the automatic channel has a strongly negative quadratic slope while the judge channel is flat (close to zero, not significant); the two channels agree on \emph{direction} only in the loosest sense and the judge channel does not yet reveal a competing positive curvature.
For \textbf{Qwen3-8B}, the signs already \textbf{invert} in point estimates ($\beta_{d^2}^{\mathrm{auto}}{=}{-}0.744$ vs.\ $\beta_{d^2}^{\mathrm{judge}}{=}{+}0.226$), but only the automatic side is significant; this is the regime where the two scoring channels start pointing in opposite directions without yet producing dual significance.
For \textbf{Qwen3-14B}, both quadratic terms become \textbf{simultaneously} significant with opposite signs ($\beta_{d^2}^{\mathrm{auto}}{=}{-}0.432$, $p{\approx}1.2{\times}10^{-4}$; $\beta_{d^2}^{\mathrm{judge}}{=}{+}0.700$, $p{\approx}5.0{\times}10^{-3}$). This is direct evidence that the two channels measure different constructs at the same model: the rubric-based judge identifies a high-conflict \emph{lift}, while the NLI$\times$perplexity composite identifies a high-conflict \emph{penalty}, on the \emph{same} T3 rows.

The proprietary \textbf{GPT-4o} arm is the strongest case (Table~\ref{tab:judge_auto_divergence}): $\beta_{d^2}^{\mathrm{auto}}{=}{-}0.618$ ($p{\approx}7.0{\times}10^{-10}$) and $\beta_{d^2}^{\mathrm{judge}}{=}{+}1.007$ ($p{\approx}1.8{\times}10^{-5}$).
Both are highly significant yet point in opposite directions, which is difficult to reconcile with a single latent ``creativity'' unless one allows the instruments to weight different stress responses.
We interpret this as systematic \textbf{construct divergence}: the automatic composite primarily penalizes \emph{semantic drift} away from the source under rising conflict, whereas the judge composite rewards \emph{expressive surprise} encoded in the rubric---both are nominally about ``creative'' rewrites, but they encode different definitions of what counts as a gain.
Construct divergence is therefore not a claim of \emph{zero} linear correlation between composites---Table~\ref{tab:validity_corr} shows moderate positive alignment---but of \textbf{different curvature} along $d_H$: row-level agreement can coexist with opposite $\beta_{d^2}$ estimates when stress reallocates sub-dimensions differently across channels.
At Qwen3-14B and on the proprietary GPT pool, both quadratic terms are significant with opposite signs on the \emph{same} T3 rows, making the split hard to attribute to sampling noise alone.

\paragraph{Surprise dimension along $d_H$.}
The rubric \emph{surprise} score (Likert mean$/5$) can be regressed with the \textbf{identical} T3-only mixed specification as $C_{\mathrm{auto}}$ and $C_{\mathrm{judge}}$: $\mathrm{surprise}{\sim} d_H{+}d_H^2{+}C(\mathrm{genre})$ with a source random intercept on the \textbf{same} $n{=}1440$ proprietary T3 discrete rows as Table~\ref{tab:judge_auto_divergence} (half GPT-4o, half GPT-4o mini).
Table~\ref{tab:surprise_curvature} shows that \textbf{surprise} carries a large positive quadratic coefficient ($\beta_{d^2}{\approx}{+}1.89$, $p{<}10^{-8}$), closely tracking $C_{\mathrm{judge}}$ ($\beta_{d^2}{\approx}{+}1.01$) and opposite $C_{\mathrm{auto}}$ ($\beta_{d^2}{\approx}{-}0.62$).
Further splitting the rubric shows that the surprise dimension's positive curvature is \textbf{almost twice} that of the headline judge composite and an order of magnitude tighter in $p$-value, so the high-conflict \emph{lift} on $C_{\mathrm{judge}}$ concentrates overwhelmingly on human-rated \emph{surprise} rather than on a uniform shift across all seven Likert legs.
At the open-source scale, the same lift now reappears on Qwen3-14B per-dimension fits: \emph{surprise} carries $\beta_{d^2}{\approx}{+}1.97$ ($p{\approx}6.6{\times}10^{-9}$) and \emph{imagery} $\beta_{d^2}{\approx}{+}2.62$ ($p{\approx}4.6{\times}10^{-11}$), again opposite the automatic channel.
In short: high conflict raises judged quality chiefly by activating perceived surprise (and imagery), which sharpens the construct-divergence story with concrete rubric sources at both proprietary and open-source scales.

\begin{table}[t]
\centering
\small
\begin{tabular}{lrr}
\toprule
Outcome (T3 pooled, $n{=}1440$) & $\beta_{d^2}$ & $p(d^2)$ \\
\midrule
Surprise$/5$ (rubric mean) & $+1.891$ & $8.1{\times}10^{-9}$ \\
$C_{\mathrm{auto}}$ & $-0.618$ & $7.0{\times}10^{-10}$ \\
$C_{\mathrm{judge}}$ & $+1.007$ & $1.8{\times}10^{-5}$ \\
\bottomrule
\end{tabular}
\caption{T3-only quadratic curvature for rubric surprise vs.\ headline composites (Gaussian mixed-effects with genre fixed effects and source random intercepts; proprietary discrete-main pool, $n{=}1440$).}
\label{tab:surprise_curvature}
\end{table}

\subsection{CrEval pairwise preference arm}

To add a third evaluation channel, we score a \textbf{matched discrete-main subset} with \textbf{CrEval}, a public pairwise creativity evaluator for text~\cite{cao2025creval,creval_repo}. The acronym overlaps with the image-domain \textbf{CREval} benchmark of Wang et al.~\cite{wang2026crevalimage}; throughout this report, \emph{CrEval} refers specifically to the \emph{text} evaluator used for pairwise rewrite comparison.

The scored slice uses one canonical rewrite per condition-target cell (\texttt{repeat\_idx}{=}0) from the discrete main arm: $n{=}1680$ rewrites in total across GPT-4o and Qwen3-4B/8B/14B. For each of the 60 source texts, all available rewrites are placed into a \textbf{round-robin tournament}: every distinct pair is judged exactly once, with response order randomized independently per comparison to suppress left--right positional bias. This yields \textbf{22,680} pairwise decisions overall, with each rewrite facing exactly 27 opponents.

Implementation follows the public CrEval release (Qwen2.5-7B-Instruct base with the CrEval-7b LoRA adapter served through a LLaMA-Factory OpenAI-compatible API)~\cite{cao2025creval,creval_repo}. The model returns one of three outcomes for each pair: response 1 preferred, response 2 preferred, or tie. We convert these into a win-rate
\[
  C_{\mathrm{creval}} = \frac{\mathrm{wins} + 0.5\,\mathrm{ties}}{n_{\mathrm{comparisons}}},
\]
which lies in $[0,1]$ by construction and is interpreted as a \textbf{pairwise preference score} rather than an absolute rubric rating.

This channel does \emph{not} replace the prespecified primary outcome $C_{\mathrm{auto}}$. Its role is triangulatory: $C_{\mathrm{auto}}$ remains the fidelity-sensitive automatic composite, $C_{\mathrm{judge}}$ remains the absolute human-like rubric channel, and $C_{\mathrm{creval}}$ adds a \textbf{calibrated pairwise preference arm} that asks which rewrite is preferred overall when two candidates are viewed head-to-head.

\subsection{CrEval relative to the absolute judge arm}

On this matched 1680-row slice, CrEval is positively aligned with the existing judge channel but is not redundant with it (Table~\ref{tab:creval_alignment}). Pearson correlation is moderate between $C_{\mathrm{creval}}$ and $C_{\mathrm{judge}}$ ($r{=}0.57$), stronger with the single GPT-4o judge stream than with the DeepSeek stream ($r{=}0.60$ vs.\ 0.49), and strongest with the rubric \emph{surprise} dimension ($r{=}0.727$). By contrast, correlation with $C_{\mathrm{auto}}$ remains weak ($r{=}0.194$). The same construct split seen in Table~\ref{tab:judge_auto_divergence} therefore survives this matched-slice check: pairwise preference clusters with human-like surprise more than with the automatic NLI$\times$perplexity channel.

CrEval also changes where the evaluation signal is sharpest. On the coarse T0--T3 manipulation, all three channels separate the identity floor from rewritten outputs; in fact, on this matched subset the absolute judge composite is slightly more condition-sensitive than the pairwise score ($\eta^2{=}0.191$ for $C_{\mathrm{judge}}$ vs.\ 0.179 for $C_{\mathrm{creval}}$; Table~\ref{tab:creval_alignment}). The gain from CrEval lies elsewhere: when the factor is \textbf{generator identity} rather than condition, $C_{\mathrm{creval}}$ explains 7.8\% of variance, versus 2.3\% for $C_{\mathrm{judge}}$ and 3.0\% for $C_{\mathrm{auto}}$ (Figure~\ref{fig:creval_channels}; Table~\ref{tab:creval_alignment}). In other words, the pairwise evaluator is not a universal replacement for the absolute judge, but it is a \textbf{sharper instrument for cross-model ranking} on the matched slice actually scored with CrEval.

\begin{table}[t]
\centering
\small
\begin{tabular}{lrrrrr}
\toprule
Condition & Creativity & Novelty & Value & NLI & Coherence \\
\midrule
T0 & 0.2317 & 0.3334 & 0.7044 & 0.7486 & 0.6602 \\
T1 & 0.3197 & 0.4176 & 0.7677 & 0.7759 & 0.7596 \\
T2 & 0.3120 & 0.4231 & 0.7408 & 0.7153 & 0.7662 \\
T3 & 0.3124 & 0.4232 & 0.7409 & 0.7144 & 0.7675 \\
\midrule
\multicolumn{6}{@{}l@{}}{\footnotesize\emph{Open-source Qwen3 generators, T3 condition only.}} \\
Qwen3-4B & 0.2785 & 0.3700 & 0.7616 & 0.8022 & 0.7211 \\
Qwen3-8B & 0.2876 & 0.3812 & 0.7635 & 0.7857 & 0.7413 \\
Qwen3-14B & 0.2967 & 0.3859 & 0.7759 & 0.8162 & 0.7355 \\
\midrule
M0 & 0.1615 & 0.4415 & 0.3879 & 0.4147 & 0.3611 \\
M1 & 0.2044 & 0.4471 & 0.4771 & 0.5652 & 0.3890 \\
M2 & 0.1908 & 0.4583 & 0.4305 & 0.5372 & 0.3238 \\
M3 & 0.2018 & 0.4013 & 0.5084 & 0.5305 & 0.4863 \\
\bottomrule
\end{tabular}
\caption{Condition means. T0--T3: proprietary discrete main arm. Qwen3 rows: T3 only on the same protocol. M0--M3: mechanism arm, fifteen independent generations per mode on top-$d$ cells.}
\label{tab:condition}
\end{table}

\section{Descriptive Results}

\subsection{Genre-Level Patterns}

Pooling discrete and continuous main arms (4560 rows) yields the following means.

\begin{table}[t]
\centering
\small
\begin{tabular}{lrrrrr}
\toprule
Genre & Creativity & Novelty & Value & NLI & Coherence \\
\midrule
Academic & 0.2838 & 0.3913 & 0.7342 & 0.6879 & 0.7805 \\
Essay & 0.3434 & 0.4466 & 0.7730 & 0.8078 & 0.7383 \\
Narrative & 0.2593 & 0.3628 & 0.7174 & 0.5547 & 0.8801 \\
Poetry & 0.3385 & 0.4520 & 0.7519 & 0.9054 & 0.5984 \\
\bottomrule
\end{tabular}
\caption{Genre means pooling discrete and continuous main arms (4560 rows, proprietary generators).}
\label{tab:genre}
\end{table}

\begin{figure*}[t]
\centering
\begin{minipage}[t]{0.49\textwidth}
\centering
\maybeincludegraphics{figures/inverted_u_creativity_auto.png}{Missing figure.}
\end{minipage}\hfill
\begin{minipage}[t]{0.49\textwidth}
\centering
\maybeincludegraphics{figures/directional_field_creativity_auto.png}{Missing figure.}
\end{minipage}
\caption{Binned mean $C_{\mathrm{auto}}$ vs.\ $d_H$ and the directional field in $(\Delta S,\Delta R)$ (proprietary pooled sample). Headline creativity trends downward with $d_H$ in pooled binned means (left); regressions still estimate negative $d_H^2$ on $C_{\mathrm{auto}}$ over the mixed sample alongside positive novelty curvature (Table~\ref{tab:quadratic}). Right: interactive $\Delta S$--$\Delta R$ structure.}
\label{fig:main}
\end{figure*}

\begin{figure*}[t]
\centering
\begin{minipage}[t]{0.49\textwidth}
\centering
\maybeincludegraphics{figures/judge_auto_validity_matrix.png}{Missing figure.}
\end{minipage}\hfill
\begin{minipage}[t]{0.49\textwidth}
\centering
\maybeincludegraphics{figures/t3_d_H_distribution_main.png}{Missing figure.}
\end{minipage}
\caption{Left: Pearson correlations among judge sub-scores and automatic component scores (experiment-wide flat table). Right: $d_H$ distribution on discrete-main T3 rows, contextualizing the H1 vertex relative to observed conflict mass.}
\label{fig:aux}
\end{figure*}

\begin{figure}[t]
\centering
\maybeincludegraphics{figures/generation_variability_sentiment_kendall.png}{Missing figure.}
\caption{Exploratory \textbf{generation variability} vs.\ conflict: eight equal-width $d_H$ bins over the pooled sample. \emph{Top:} mean absolute sentiment displacement ($[0,2]$). \emph{Bottom:} mean sentence-order preservation score ($[0,1]$, higher $\approx$ closer to source order).}
\label{fig:variability}
\end{figure}

\section{Conflict Intensity}

\subsection{Regression specification}

The pooled quadratic models regress each outcome on linear and quadratic conflict terms plus categorical genre and experimental condition:
\[
  y \sim d_H + d_H^2 + \sum_g \mathbf{1}_{\mathrm{genre}=g} + \sum_c \mathbf{1}_{\mathrm{condition}=c}.
\]
Observations are rewrite-level rows; \textbf{random intercepts group by source text}, reflecting repeated generations from the same corpus item without nesting random effects in persona identity. Condition absorbs arm contrasts (T0--T3, mechanism modes, multihop structure as coded); persona distance enters through $d_H$ and directional regressors rather than as an extra random classification. Mixed linear models are used when variance-component estimation is numerically stable; otherwise the same fixed-effects structure is estimated by OLS. In sparse designs the estimated source-level variance may collapse toward zero, so $\beta_{d^2}$ should be read as a pooled association that partly marginalizes within-source correlation rather than as a nested hierarchical causal estimand.
Directional regressions swap $(d_H,d_H^2)$ for $(\Delta S,\Delta R)$ plus squares, product, and Euclidean $|\Delta_H|$, retaining the same genre/condition fixed effects and source-level intercept grouping.

\subsection{Pooled Quadratic Results}

Mixed-effects fits surface the dimensional split: \textbf{negative} $d_H^2$ on headline creativity, value, and coherence versus \textbf{positive} $d_H^2$ on novelty (combined), over $n{=}4620$ rows with valid $d_H$ and finite outcomes on each metric (Table~\ref{tab:quadratic}; proprietary pooled sample).
Together with the T3-only vertex $d^\ast{\approx}0.47$ (Table~\ref{tab:h1}), one can read a medium-conflict regime as comparatively favourable \emph{within that subset}, in line with classic arousal--performance tradeoffs~\cite{yerkes}; Figure~\ref{fig:main} nonetheless shows pooled binned $C_{\mathrm{auto}}$ declining with $d_H$, so the vertex should not be equated with a visible global inverted-U in the aggregate bins.

\paragraph{Scale dependence of the inverted-U.}
Table~\ref{tab:scale_regression} fixes the same quadratic specification on automatic creativity ($C_{\mathrm{auto}}\sim d_H+d_H^2+C(\mathrm{genre})$) with \textbf{source-level random intercepts}, fit \textbf{separately per metrics file} on \textbf{T3 discrete-target rows only}.
Each open-source Qwen checkpoint file contributes $n{=}720$ T3 rows (one generator model per file).
The proprietary slice pools \textbf{all} T3 discrete-target rows for the two OpenAI generator assignments in a \emph{single} mixed-effects fit ($n{=}1440$: half GPT-4o, half GPT-4o mini), matching Table~\ref{tab:judge_auto_divergence}. This is not ``half the corpus'' and not a discrete-vs-continuous split: continuous-target T3 rows for these generators are not included in that proprietary discrete pool.
Optimisers use L-BFGS-B by default; Powell is used when the Hessian is singular.

\begin{table}[t]
\centering
\small
\begin{tabular}{lrrrl}
\toprule
Generator (T3) & $\beta_{d^2}$ & $p(d^2)$ & $d^\ast$ & Sig.\ ($p{<}0.05$) \\
\midrule
Qwen3-4B & $-0.6408$ & $1.74{\times}10^{-9}$ & 0.5305 & Yes \\
Qwen3-8B & $-0.7440$ & $2.98{\times}10^{-8}$ & 0.4611 & Yes \\
Qwen3-14B & $-0.4323$ & $1.21{\times}10^{-4}$ & 0.4552 & Yes \\
GPT-4o (pooled) & $-0.6181$ & $6.95{\times}10^{-10}$ & 0.4954 & Yes \\
\bottomrule
\end{tabular}
\caption{Quadratic curvature for $C_{\mathrm{auto}}$ vs.\ $d_H$ by generator (T3-only discrete rows; mixed-effects with source intercepts). Qwen checkpoints: $n{=}720$ each. GPT-4o (pooled): $n{=}1440$ (half GPT-4o, half GPT-4o mini) in one fit, matching Table~\ref{tab:judge_auto_divergence}.}
\label{tab:scale_regression}
\end{table}

\begin{table}[t]
\centering
\small
\begin{tabular}{lrrr}
\toprule
Metric & $n$ & $\beta_{d^2}$ & $p(d^2)$ \\
\midrule
Creativity auto & 4620 & $-0.1111$ & $4.58{\times}10^{-7}$ \\
Creativity geom & 4620 & $-0.1479$ & $2.76{\times}10^{-8}$ \\
Value auto & 4620 & $-0.3901$ & $7.22{\times}10^{-17}$ \\
Coherence auto & 4620 & $-0.3266$ & $5.37{\times}10^{-11}$ \\
Novelty (combined) & 4620 & $0.0908$ & $1.90{\times}10^{-6}$ \\
\bottomrule
\end{tabular}
\caption{Quadratic terms from pooled mixed-effects regressions on headline outcomes and components (proprietary pooled sample). Singular random-effect covariance warnings can occur at numerical boundaries.}
\label{tab:quadratic}
\end{table}

\subsection{Pre-Registered H1 Check}

The fitted vertex $d^\ast{\approx}0.47$ on the proprietary T3 mass operationalizes \textbf{medium} conflict (empirical q03--q97 shown); it lies outside the narrower pre-registered vertex band yet inside the sample-mass diagnostic, so claims should reference both the band logic and the descriptive optimum. This interior optimum is estimated on T3 rows only and coexists with monotonic-looking pooled binned means vs.\ $d_H$ (Figure~\ref{fig:main}). All three Qwen3 vertices in Table~\ref{tab:scale_regression} fall inside the per-checkpoint empirical q03--q97 mass (e.g.\ 4B: $[0.295,0.797]$, vertex 0.531; 8B: same support, vertex 0.461; 14B: vertex 0.455).

\begin{table}[t]
\centering
\small
\begin{tabular}{lr}
\toprule
H1 diagnostic & Value \\
\midrule
T3 main rows & 1440 \\
$\beta_d$ & 0.3265 \\
$\beta_{d^2}$ & $-0.3481$ \\
$p(d^2)$ & $6.51{\times}10^{-4}$ \\
Fitted vertex $d$ & 0.4689 \\
Empirical $d$ q03 / q97 & 0.2948 / 0.7971 \\
Vertex inside pre-reg band & False \\
Vertex inside sample mass & True \\
Pooled main $p(d^2)$ & $4.58{\times}10^{-7}$ \\
\bottomrule
\end{tabular}
\caption{Pre-registered H1 diagnostic for automatic creativity (proprietary discrete-main T3).}
\label{tab:h1}
\end{table}

\paragraph{Per-checkpoint H1 (Qwen3-4B / 8B / 14B).}
The \texttt{qwen\_h\_full\_analysis} script reruns the H1 OLS pieces on each open-source checkpoint at $n{=}720$ T3 rows (Table~\ref{tab:qwen_h1}). The MixedLM estimates in Table~\ref{tab:scale_regression} and the OLS-with-genre-FE estimates here use the same quadratic form but can differ slightly in $p$-values because of the random-intercept grouping; both agree that all three checkpoints place a vertex inside the empirical $d_H$ mass.

\begin{table}[t]
\centering
\small
\begin{tabular}{lrrr}
\toprule
H1 diagnostic & Qwen3-4B & Qwen3-8B & Qwen3-14B \\
\midrule
$n$ T3 & 720 & 720 & 720 \\
$\beta_d$ (OLS) & 0.5338 & 0.5615 & 0.2242 \\
$\beta_{d^2}$ (OLS) & $-0.5023$ & $-0.6259$ & $-0.2721$ \\
$p(d^2)$ & $1.33{\times}10^{-4}$ & $4.63{\times}10^{-5}$ & $4.92{\times}10^{-2}$ \\
Fitted vertex $d$ & 0.5314 & 0.4485 & 0.4120 \\
Empirical $d$ q03 / q97 & 0.2948 / 0.7971 & 0.2948 / 0.7971 & 0.2948 / 0.7971 \\
Vertex inside sample mass & True & True & True \\
\bottomrule
\end{tabular}
\caption{Per-checkpoint H1 for the open-source Qwen3 replication (T3 only, OLS with genre fixed effects). MixedLM $\beta_{d^2}$ for 14B remains significant at $p{\approx}1.2{\times}10^{-4}$ (Table~\ref{tab:scale_regression}); the OLS $p(d^2){=}0.049$ here is borderline at the conventional threshold.}
\label{tab:qwen_h1}
\end{table}

\begin{figure}[t]
\centering
\maybeincludegraphics{figures/scale_inverted_u.png}{Missing figure.}
\caption{T3-only genre-averaged OLS predictions of $C_{\mathrm{auto}}$ vs.\ $d_H$ for Qwen3 checkpoints (4B/8B/14B). Vertical dotted lines mark fitted vertices for 4B/8B/14B; shaded bands sketch approximate hybrid intervals around each vertex (the per-checkpoint band boundaries shown in the figure are read off the curve crossings used by \texttt{plot\_scale\_inverted\_u.py} and are illustrative rather than inferential).}
\label{fig:scale_inverted_u}
\end{figure}

\begin{figure}[t]
\centering
\maybeincludegraphics{figures/scale_inverted_u_judge.png}{Missing figure.}
\caption{Same axes and curve construction as Figure~\ref{fig:scale_inverted_u}, but with mean predicted $C_{\mathrm{judge}}$ (genre-averaged OLS) on the vertical axis. Contrasts the rubric-based composite against the automatic $C_{\mathrm{auto}}$ trajectories on identical Qwen3 T3 rows; the upward-bending 14B trace parallels the GPT-4o judge-side pattern.}
\label{fig:scale_inverted_u_judge}
\end{figure}

\subsection{Conflict direction}

\begin{table}[t]
\centering
\small
\begin{tabular}{lrr}
\toprule
Term & Coefficient & $p$ \\
\midrule
$\Delta S$ & 0.0164 & $3.57{\times}10^{-43}$ \\
$\Delta R$ & 0.0028 & 0.0127 \\
$\Delta S^2$ & $-0.0018$ & 0.607 \\
$\Delta R^2$ & $-0.0143$ & $3.67{\times}10^{-3}$ \\
$\Delta S\Delta R$ & 0.0075 & $1.93{\times}10^{-11}$ \\
$|\Delta|$ & 0.0190 & 0.0343 \\
\bottomrule
\end{tabular}
\caption{Directional regression for automatic creativity (proprietary pooled sample).}
\label{tab:direction}
\end{table}

The same directional structure reappears when we restrict to the \textbf{open-source Qwen3-4B/8B/14B} T3 slices: the interaction term $\Delta S{\times}\Delta R$ remains highly significant on all three checkpoints ($p{\approx}3.5{\times}10^{-5}$ on 4B, $p{\approx}1.4{\times}10^{-3}$ on 8B, $p{<}10^{-6}$ on 14B), and the negative $\Delta S^2$, $\Delta R^2$ curvature pair persists on 4B/8B alongside a positive $|\Delta|$ coefficient---matching the pooled GPT-family pattern in Table~\ref{tab:direction}. On 14B, $\Delta R$ alone is strongly linear ($p{\approx}9.1{\times}10^{-5}$) while $\Delta S^2$ and $\Delta R^2$ are weaker, suggesting that register shifts dominate the directional channel at the largest open-source scale.

\subsection{Causal contrasts}

Bucketed contrasts pool \textbf{both} main metric files (hence larger per-bucket $n$ than a discrete-only table; proprietary pooled sample).

Across all three conflict buckets, mean $C_{\mathrm{auto}}$ at T3 lies below both T1 and T2 (Table~\ref{tab:causal}), and both T3--T1 and T3--T2 contrasts are negative or near zero. At this pooled scale, \textbf{directed cross-persona targeting (T3) does not outperform} nearest-neighbour (T1) or random persona assignment (T2) on automatic creativity; if anything, it trails both---a pattern that should temper expectations that ``designed'' cross-persona moves automatically beat placebo-style persona swaps.

\begin{table}[t]
\centering
\footnotesize
\setlength{\tabcolsep}{4pt}
\begin{tabular}{@{}lrrrrr@{}}
\toprule
Bucket & $n_{T1}$ & $n_{T2}$ & $n_{T3}$ & T3$-$T1 & T3$-$T2 \\
\midrule
Low  & 228 &  68 & 1004 & $-0.0023$ & $-0.0009$ \\
Mid  & 468 & 240 &  841 & $-0.0165$ & $-0.0061$ \\
High &  24 & 412 & 1035 & $-0.0523$ & $-0.0078$ \\
\bottomrule
\end{tabular}
\par\medskip
\begin{tabular}{@{}lrrr@{}}
\toprule
Bucket & $\bar{C}_{T1}$ & $\bar{C}_{T2}$ & $\bar{C}_{T3}$ \\
\midrule
Low  & 0.3127 & 0.3113 & 0.3104 \\
Mid  & 0.3211 & 0.3107 & 0.3046 \\
High & 0.3591 & 0.3146 & 0.3068 \\
\bottomrule
\end{tabular}
\caption{Bucketed causal contrasts on $C_{\mathrm{auto}}$ (proprietary pooled sample). Top: cell counts and mean contrasts; bottom: condition means.}
\label{tab:causal}
\end{table}

On the \textbf{Qwen3} replications, the H6 mid-bin contrast $C_{\mathrm{auto}}^{T3}{-}C_{\mathrm{auto}}^{T2}$ is again \emph{near zero} on 4B ($-0.0016$) and 8B ($-0.0041$), but turns \textbf{slightly positive} on 14B ($+0.0039$), a small-scale reversal that does not overturn the proprietary pattern but hints that the largest checkpoint may occasionally edge out the random-persona placebo in the mid-$d_H$ band. For H5, JSD-curvature on $C_{\mathrm{auto}}$ is non-significant on 4B ($p{\approx}0.41$) and negatively curved on 8B ($\beta_{\mathrm{JSD}^2}{\approx}{-}7.07$, $p{\approx}0.017$), and again non-significant on 14B ($p{\approx}0.88$). Directed T3 persona choice does not deliver incremental gains over the T2 placebo within fixed $d_H$ bins on any Qwen scale at conventional thresholds, paralleling the proprietary pattern. We continue to report this as a substantive empirical finding rather than as a methodological defect.

\paragraph{Structural tail check (appendix).}
The automatic-channel deterioration under high conflict is already conveyed by $C_{\mathrm{auto}}$, NLI, and coherence slopes in the main tables.
For readers who want a \textbf{binary} articulation of the same stress axis, Appendix~\ref{app:collapse} fits a clustered binomial GEE on a joint low-NLI-and-low-coherence indicator (sample-relative $q{=}0.2$ cut-offs on the same $n{=}1440$ proprietary T3 pool); marginal fitted risk rises steeply at high $d_H$.
We keep that material in an appendix because it introduces GEE machinery and sample-defined thresholds that sit awkwardly with the main mixed-effects narrative, and because the qualitative message largely overlaps what is already visible in the automatic scores.

\subsection{CrEval on the matched discrete-main slice}

Figure~\ref{fig:creval_channels}(a) places the new pairwise channel beside $C_{\mathrm{judge}}$ and $C_{\mathrm{auto}}$ on the same 1680 matched rows. The main condition result is unambiguous: the identity floor T0 is far below every rewriting condition under all three metrics. For $C_{\mathrm{creval}}$, mean win-rate rises from 0.2236 at T0 to 0.5069/0.5539/0.5539 for T1/T2/T3; the T3--T0 contrast carries Cohen's $d{=}1.375$, close to the matched-slice judge effect ($d{=}1.419$) and substantially larger than the automatic-channel contrast ($d{=}0.922$). Thus CrEval confirms the core finding that persona rewriting is preferred to identity preservation; its distinctive value is not the existence of a condition effect, but the way it reweights \emph{which} rewrites are preferred within the rewritten pool.

That reweighting is most visible on the generator axis. Figure~\ref{fig:creval_channels}(b) shows the largest gap between GPT-4o and the three Qwen checkpoints under $C_{\mathrm{creval}}$: mean scores are 0.629 for GPT-4o versus 0.441--0.466 for Qwen3, and the associated model-factor ANOVA reaches $\eta^2{=}0.078$ compared with 0.023 for $C_{\mathrm{judge}}$ and 0.030 for $C_{\mathrm{auto}}$ (Table~\ref{tab:creval_alignment}). On this matched subset, the pairwise evaluator is therefore the \textbf{most discriminative of generator capacity}, even though the absolute judge remains slightly stronger on the coarse T0--T3 manipulation.

The high-conflict curvature check leads to the same conclusion. On T3 rows only ($n{=}240$ per generator in the CrEval slice), $C_{\mathrm{creval}}$ detects positive quadratic curvature for Qwen3-14B and Qwen3-8B ($\beta_{d^2}{=}+2.022$, $p{=}0.010$; $\beta_{d^2}{=}+1.615$, $p{=}0.031$), whereas the same-slice $C_{\mathrm{judge}}$ coefficients stay positive but non-significant (Table~\ref{tab:creval_t3_quadratic}). GPT-4o behaves differently: $C_{\mathrm{auto}}$ still penalizes higher conflict on the matched T3 slice ($\beta_{d^2}{=}{-}0.518$, $p{=}0.027$), but neither $C_{\mathrm{judge}}$ nor $C_{\mathrm{creval}}$ shows reliable curvature, consistent with a pairwise \textbf{quality ceiling} once generator quality is already high.

Descriptively, CrEval also shifts the directional emphasis of the story. Within T3, academic rows have the highest mean pairwise score (0.602) and poetry the lowest (0.515), reversing the novelty-heavy ordering of the automatic channel; adventurous target personas dominate conservative ones by roughly 0.35--0.45 win-rate units. In that sense, the pairwise preference arm weights \textbf{register lift and perceived surprise} more heavily than the automatic channel, complementing the directional evidence in Table~\ref{tab:direction} rather than duplicating it.

\section{Mechanism and Multihop Branches}

\subsection{Mechanism Experiment}

Mechanism generations restrict to the fifteen highest-$d$ discrete-main cells and cross four informational modes (M0--M3); each mode has fifteen rows (Table~\ref{tab:condition}).

\paragraph{Which instructional mechanism helps most?}
Pooling those sixty rows, mean automatic creativity orders \textbf{M1} (direct rewrite, highest mean $C_{\mathrm{auto}}$) ${>}$ \textbf{M3} (constraint-forward rewrite) ${>}$ \textbf{M2} (two-stage) ${>}$ \textbf{M0} (control). This addresses the project-book theme that different conflict-\emph{instruction} bundles yield different creative profiles: unconstrained cross-persona rewriting (M1) best preserves $C_{\mathrm{auto}}$ here, tightly constrained rewriting (M3) is second, explicit decomposition (M2) sits mid-pack, and the passive control (M0) is lowest---all on the same high-$d$ cells.
With $n{=}15$ per mode this ordering is \textbf{descriptive}; formal mechanism-only inference would warrant dedicated models with pair-level blocking.

Table~\ref{tab:mechanism} restricts to rows with finite $d_H$ and outcome, takes the upper third of $d_H$ \emph{globally}, then aggregates by mode. The M0--M3 lines describe mechanism modes inside that high-$d$ slice; the joint-mode line summarizes continuous-main rows under that label in the same slice ($n{=}1470$), not the sixty-row mechanism-only sample.

\begin{table}[t]
\centering
\small
\begin{tabular}{lrrr}
\toprule
Mode & Mean creativity & Std. & Count \\
\midrule
M0 & 0.1615 & 0.0701 & 15 \\
M1 & 0.2044 & 0.0793 & 15 \\
M2 & 0.1908 & 0.0859 & 15 \\
M3 & 0.2018 & 0.0674 & 15 \\
Joint & 0.3099 & 0.0765 & 1470 \\
\bottomrule
\end{tabular}
\caption{Mechanism-related summary in the global high-$d$ slice.}
\label{tab:mechanism}
\end{table}

\subsection{Multihop Prediction Check}

\begin{table}[t]
\centering
\small
\begin{tabular}{llrrr}
\toprule
Genre & Bin & Multihop & Direct T3 & Delta \\
\midrule
Academic & Low & 0.3244 & 0.2941 & 0.0303 \\
Academic & Mid & 0.2134 & 0.2830 & $-0.0696$ \\
Academic & High & 0.2553 & 0.2736 & $-0.0182$ \\
Essay & Low & 0.3968 & 0.3412 & 0.0557 \\
Essay & Mid & 0.3948 & 0.3527 & 0.0421 \\
Essay & High & 0.4142 & 0.3429 & 0.0713 \\
Narrative & Low & 0.2695 & 0.2706 & $-0.0011$ \\
Narrative & Mid & 0.2781 & 0.2515 & 0.0266 \\
Narrative & High & 0.2733 & 0.2619 & 0.0115 \\
Poetry & Low & 0.3333 & 0.3298 & 0.0035 \\
Poetry & Mid & 0.4294 & 0.3404 & 0.0890 \\
Poetry & High & 0.3704 & 0.3512 & 0.0192 \\
\bottomrule
\end{tabular}
\caption{Multihop vs.\ direct T3 mean automatic creativity by genre and $d$ bin.}
\label{tab:multihop}
\end{table}

Multihop effects are \textbf{genre-dependent}. Essay rows show multihop strictly above direct T3 in every $d$ bin (all positive deltas), suggesting smoother persona paths help that register here. Academic rows reverse under mid conflict (delta ${\approx}{-}0.07$), while low/high bins are mixed or modest---so gradual hops are not universally preferable and should be evaluated per genre and conflict band.

\section{Limitations}

\begin{itemize}\itemsep2pt
\item For smaller checkpoints (especially Qwen3-4B), \textbf{genre-averaged OLS} curves in Figure~\ref{fig:scale_inverted_u} are illustrative; their nominal standard errors do not incorporate the same hierarchical structure as the MixedLM rows in Table~\ref{tab:scale_regression} and can be \textbf{anti-conservative} when read as inferential objects rather than as sketches of shape.
\item The ``\textbf{capacity threshold}'' language for Qwen3-14B (judge curvature turning significantly positive while automatic curvature stays significantly negative on the \emph{same} rows) is a \textbf{post hoc descriptive} reading of two fitted slopes at one model size, not an independently identified causal capacity parameter.
\item Appendix~\ref{app:collapse} reports a binomial GEE on sample-defined collapse cut-offs (not a calibrated GLMM); treat as exploratory triangulation, not as an absolute failure probability.
\item H1 OLS $p(d^2)$ for Qwen3-14B is borderline at the conventional $0.05$ threshold ($p{\approx}0.049$) even though the MixedLM fit in Table~\ref{tab:scale_regression} is highly significant; readers should treat the vertex location as robust but the OLS $p$-value as estimator-sensitive.
\item Space-L reference text uses \textbf{GPT-4o / GPT-4o mini} only; H7 applies the cached L geometry to Qwen3 rewrites (Figure~\ref{fig:space_l_validation}).
\item The CrEval analysis is restricted to the discrete-main \texttt{repeat\_idx}{=}0 slice ($n{=}1680$ rows, 960 T3 rows), not the full 4716-row proprietary pooled sample nor the full 3420-row Qwen replication; its coefficients should therefore be read as a \textbf{matched-slice triangulation}, not as a replacement for the main pooled estimands.
\item CrEval win-rates are tournament-relative: each rewrite is scored against the observed pool of same-source opponents, mixing generators of different capacity inside one competition graph. This improves pairwise calibration but makes absolute levels \textbf{sample-dependent}.
\item Per-generator CrEval regressions on the matched T3 slice use OLS over $n{=}240$ rows per model because mixed-effects fits are numerically unstable in this narrow subset; Table~\ref{tab:creval_t3_quadratic} should therefore be read as a descriptive curvature check rather than as a new primary inferential table.
\end{itemize}

\section{Conclusion}

Cross-persona conflict behaves as \textbf{structural stress}: it does \emph{not} uniformly lift creativity; instead it reallocates mass across automatic subscores---novelty curvature turns positive while value and coherence turn sharply negative (Table~\ref{tab:quadratic}). The discrete-main T3-only quadratic places a vertex near $d^\ast{\approx}0.47$ (Table~\ref{tab:h1}), while pooled binned means of $C_{\mathrm{auto}}$ vs.\ $d_H$ trend downward at scale (Figure~\ref{fig:main}); together these caution against equating ``medium conflict'' with a visibly inverted-U aggregate trajectory for headline creativity.
Exploratory variability traces (Figure~\ref{fig:variability}) show sentiment displacement rising into mid-to-high conflict bins and the sentence-order score dipping through mid bins, consistent with more aggressive reordering co-occurring with the value/coherence erosion channel.
Directionally, $\Delta S$, $\Delta R^2$, and $\Delta S\Delta R$ remain central (Table~\ref{tab:direction}). On mechanism prompts over shared high-$d$ cells, mean $C_{\mathrm{auto}}$ orders M1${>}$M3${>}$M2${>}$M0 (Table~\ref{tab:condition}), a descriptive answer to ``which bundle works best'' within this arm.
The open-source \textbf{Qwen3-4B/8B/14B} runs reproduce the same qualitative story at three scales: significant negative quadratic curvature on $C_{\mathrm{auto}}$ vs.\ $d_H$ with fitted vertices near $d^\ast{\approx}0.46$--$0.53$ (Table~\ref{tab:scale_regression}; Figure~\ref{fig:scale_inverted_u}; judge-channel analogue Figure~\ref{fig:scale_inverted_u_judge}).
Space-L Procrustes alignment and H7 $d_L$ checks confirm the same negative $\beta_{d^2}$ story (Figure~\ref{fig:space_l_validation}).
At Qwen3-14B and on the proprietary GPT pool, $C_{\mathrm{judge}}$ bends upward along $d_H$ while $C_{\mathrm{auto}}$ bends downward on the same rows (Table~\ref{tab:judge_auto_divergence}), even though both composites correlate positively at the row level (Table~\ref{tab:validity_corr}).
Treat all claims as conditional on branch (main vs.\ mechanism vs.\ multihop), sampling mode, and the narrow mechanism sample size.

For practitioners, the choice of evaluation metric matters as much as the choice of conflict distance. On the automatic channel ($C_{\mathrm{auto}}$, which penalizes semantic drift), medium conflict near $d^\ast{\approx}0.46$--$0.53$ is comparatively favourable and high conflict degrades quality---making it the right metric for fidelity-sensitive tasks such as academic editing or translation. On the judge channel ($C_{\mathrm{judge}}$, which rewards human-perceived surprise), the curvature reverses: the proprietary GPT pool shows $\beta_{d^2}^{\mathrm{judge}}{=}{+}1.007$ ($p{<}10^{-4}$) and Qwen3-14B shows $\beta_{d^2}^{\mathrm{judge}}{=}{+}0.700$ ($p{\approx}5{\times}10^{-3}$), so higher conflict can yield more positively judged rewrites under rubric-based scoring at sufficient model scale. The two channels are not interchangeable; downstream design should select the metric that matches the intended use case before tuning conflict distance.

A third channel, $C_{\mathrm{creval}}$, clarifies where the judge-side signal is genuinely preference-like. On the matched discrete-main slice it remains positively aligned with $C_{\mathrm{judge}}$ ($r{=}0.57$) and especially with rubric surprise ($r{=}0.727$), but it separates generator capacity more strongly than either the absolute judge or the automatic composite (model $\eta^2{=}0.078$ vs.\ 0.023/0.030; Figure~\ref{fig:creval_channels}, Table~\ref{tab:creval_alignment}). On the same matched T3 rows, positive high-conflict curvature becomes detectable for Qwen3-8B/14B under $C_{\mathrm{creval}}$ even when the paired $C_{\mathrm{judge}}$ slice remains non-significant (Table~\ref{tab:creval_t3_quadratic}), while GPT-4o stays near a pairwise quality ceiling. This makes CrEval useful not as a wholesale replacement for the absolute rubric, but as a calibrated \textbf{pairwise complement}: $C_{\mathrm{auto}}$ remains the conservative metric for fidelity-sensitive tasks, $C_{\mathrm{judge}}$ supplies interpretable sub-dimensions, and $C_{\mathrm{creval}}$ provides the sharpest ranking signal when the question is which rewrite is preferred overall.

\begin{figure*}[t]
\centering
\maybeincludegraphics{figures/creval_channel_comparison.png}{Matched-slice comparison of CrEval, absolute judge, and automatic creativity. Left: condition means over T0--T3. Right: generator means over GPT-4o and Qwen3 checkpoints.}
\caption{Matched-slice comparison of $C_{\mathrm{creval}}$, $C_{\mathrm{judge}}$, and $C_{\mathrm{auto}}$ on the discrete-main \texttt{repeat\_idx}{=}0 rows scored with CrEval ($n{=}1680$). Left: all three channels separate T0 from rewriting conditions. Right: $C_{\mathrm{creval}}$ yields the strongest cross-model separation, especially between GPT-4o and the Qwen3 checkpoints.}
\label{fig:creval_channels}
\end{figure*}

\begin{table}[t]
\centering
\small
\begin{tabular}{llr}
\toprule
Statistic & Quantity & Value \\
\midrule
Alignment & $r(C_{\mathrm{creval}}, C_{\mathrm{judge}})$ & 0.570 \\
Alignment & $r(C_{\mathrm{creval}}, C_{\mathrm{auto}})$ & 0.194 \\
Alignment & $r(C_{\mathrm{creval}}, \mathrm{surprise}_{\mathrm{judge}})$ & 0.727 \\
Alignment & $r(C_{\mathrm{creval}}, \mathrm{GPT\mbox{-}4o\ judge})$ & 0.601 \\
Alignment & $r(C_{\mathrm{creval}}, \mathrm{DeepSeek\ judge})$ & 0.487 \\
Condition ANOVA & $\eta^2(C_{\mathrm{creval}})$ & 0.179 \\
Condition ANOVA & $\eta^2(C_{\mathrm{judge}})$ & 0.191 \\
Condition ANOVA & $\eta^2(C_{\mathrm{auto}})$ & 0.098 \\
Model ANOVA & $\eta^2(C_{\mathrm{creval}})$ & 0.078 \\
Model ANOVA & $\eta^2(C_{\mathrm{judge}})$ & 0.023 \\
Model ANOVA & $\eta^2(C_{\mathrm{auto}})$ & 0.030 \\
\bottomrule
\end{tabular}
\caption{Matched-slice alignment and discrimination statistics for the CrEval block. Pearson correlations use all available scored rows ($n{=}1680$ except DeepSeek, $n{=}1258$); ANOVA effect sizes use the same matched rows as Figure~\ref{fig:creval_channels}. CrEval is not uniformly dominant on every criterion, but it is the most discriminative channel for \emph{generator identity}.}
\label{tab:creval_alignment}
\end{table}

\begin{table}[t]
\centering
\small
\begin{tabular}{lrrrrrr}
\toprule
Generator & $\beta_{d^2}^{\mathrm{creval}}$ & $p$ & $\beta_{d^2}^{\mathrm{judge}}$ & $p$ & $\beta_{d^2}^{\mathrm{auto}}$ & $p$ \\
\midrule
GPT-4o & $-0.664$ & 0.383 & $-0.649$ & 0.109 & $-0.518$ & 0.027 \\
Qwen3-14B & $+2.022$ & 0.010 & $+0.407$ & 0.352 & $-0.129$ & 0.520 \\
Qwen3-8B & $+1.615$ & 0.031 & $+0.255$ & 0.597 & $-0.299$ & 0.178 \\
Qwen3-4B & $+1.096$ & 0.131 & $+0.382$ & 0.395 & $+0.030$ & 0.878 \\
\bottomrule
\end{tabular}
\caption{T3-only quadratic curvature on the matched CrEval slice ($n{=}240$ rows per generator; OLS with $y{\sim}d_H{+}d_H^2$). Positive $\beta_{d^2}^{\mathrm{creval}}$ becomes detectable for Qwen3-14B and Qwen3-8B even when the same-slice absolute judge remains non-significant, while GPT-4o stays near a pairwise ceiling.}
\label{tab:creval_t3_quadratic}
\end{table}

\appendix
% Appendix: boxed ``system / user'' prompt cards (cf. ranking-prompt appendix style).
\begingroup\sloppy

\section{Appendix A: Prompt templates (boxed layout)}
\label{app:prompts}

\noindent\textbf{Implementation note.}
Full discrete persona bodies, coordinate-scorer definitions, and paraphrase variants live in the project's configuration; generation, mechanism, judging, and experiment orchestration follow the accompanying open-source implementation in the repository.

\subsection{A.1 Main arms: T0--T3 rewrites}

\subsubsection{T0 --- Identity (measurement floor)}
\begin{DefnBox}{Pipeline note}
No LLM call: the stored row is the source text verbatim with $d_H{=}0$. No system/user messages.
\end{DefnBox}

\subsubsection{T1--T3 discrete --- Joint generator (\texttt{rewrite\_joint})}
\begin{PromptSystemBox}{Rewrite / Joint --- System input}
\footnotesize
\texttt{<persona.system\_prompt>} from the persona configuration (primary or paraphrase variant), plus the length line \texttt{Target length: L--H words.}
\end{PromptSystemBox}
\begin{PromptUserBox}{Rewrite / Joint --- User input}
\LstPrompt{rewrite_user_template.txt}
\end{PromptUserBox}

\subsubsection{T3 continuous target --- Synthetic persona header}
Continuous sampling builds a lightweight synthetic persona header; numeric $S,R$ are filled per draw. The user message is the same joint wrapper as above.
\begin{PromptSystemBox}{Rewrite / Continuous T3 --- System input (pattern)}
\bgroup\color{white}\footnotesize\ttfamily\raggedright
Adopt this continuous style target in Space-H: S=+0.000, R=+0.000.\par
Lean toward rational and adventurous expression proportionally to magnitudes.\par
\egroup
\end{PromptSystemBox}
\begin{PromptUserBox}{Rewrite / Continuous T3 --- User input}
\LstPrompt{rewrite_user_template.txt}
\end{PromptUserBox}

\subsection{A.2 Mechanism arms: M0--M3}

\subsubsection{M1 --- Joint (same as main joint)}
One-shot \texttt{rewrite\_joint}: system $=$ persona $+$ length; user $=$ shared wrapper (Section~A.1).

\subsubsection{M2 --- Sequential paraphrase then persona}
\begin{PromptSystemBox}{Mechanism M2 / Stage 1 --- System input}
\bgroup\color{white}\footnotesize\ttfamily\raggedright
You are a faithful paraphraser. You restate the source's propositional content in neutral, unadorned English prose. Do not add, remove, or alter any commitment of the source. Keep the length within 90--110\% of the source.\par
\egroup
\end{PromptSystemBox}
\begin{PromptUserBox}{Mechanism M2 / Stage 1 --- User input}
\LstPrompt{paraphrase_user.txt}
\end{PromptUserBox}
\begin{PromptSystemBox}{Mechanism M2 / Stage 2 --- System input}
\footnotesize
\texttt{<persona.system\_prompt>} $+$ length constraint (same as joint).
\end{PromptSystemBox}
\begin{PromptUserBox}{Mechanism M2 / Stage 2 --- User input}
\LstPrompt{m2_stage2_user_template.txt}
\end{PromptUserBox}

\subsubsection{M3 --- Checklist-constrained joint}
\begin{PromptSystemBox}{Mechanism M3 / Checklist extraction --- System input}
\bgroup\color{white}\footnotesize\ttfamily\raggedright
You are a careful reading-comprehension assistant. From the passage below, extract a numbered checklist of its propositional commitments (factual claims, causal links, and concrete imagery) that any faithful rewrite must preserve. Be concise.\par
\egroup
\end{PromptSystemBox}
\begin{PromptUserBox}{Mechanism M3 / Checklist extraction --- User input}
\LstPrompt{checklist_user.txt}
\end{PromptUserBox}
\begin{PromptSystemBox}{Mechanism M3 / Persona + checklist --- System input}
\footnotesize
\texttt{<persona.system\_prompt>} $+$ length, followed by the injected checklist banner (trimmed to a word budget).
\end{PromptSystemBox}
\begin{PromptUserBox}{Mechanism M3 --- Checklist banner appended to system block}
\LstPrompt{m3_inject_header.txt}
\end{PromptUserBox}
\begin{PromptUserBox}{Mechanism M3 --- User message (same as joint)}
\LstPrompt{rewrite_user_template.txt}
\end{PromptUserBox}

\subsubsection{M0 --- Checklist control (token-pressure placebo)}
Checklist is still extracted for auditing, but the text appended to the persona system block is a length-matched placeholder list (control path).
\begin{PromptUserBox}{Mechanism M0 --- Control banner (system block append)}
\LstPrompt{m0_control_header.txt}
\end{PromptUserBox}
\begin{PromptUserBox}{Mechanism M0 --- User message (same as joint)}
\LstPrompt{rewrite_user_template.txt}
\end{PromptUserBox}

\subsection{A.3 Judge prompts (absolute and pairwise)}

\subsubsection{Absolute rubric (seven Likert dimensions)}
\begin{PromptSystemBox}{Judge / Absolute --- System input}
\texttt{You are a strict JSON-only literary judge.}
\end{PromptSystemBox}
\begin{PromptUserBox}{Judge / Absolute --- User input}
\LstPrompt{judge_absolute_user.txt}
\end{PromptUserBox}
\begin{DefnBox}{JSON parsing and aggregation}
Judge completions are parsed as JSON. Each Likert dimension is clipped to $[1,5]$; per-dimension scores are averaged across judges with valid JSON. A generous token ceiling is used for JSON-shaped outputs.
\end{DefnBox}

\subsubsection{Pairwise creative comparison}
\begin{PromptSystemBox}{Judge / Pairwise --- System input}
\texttt{You are a strict JSON-only literary judge.}
\end{PromptSystemBox}
\begin{PromptUserBox}{Judge / Pairwise --- User input}
\LstPrompt{judge_pairwise_user.txt}
\end{PromptUserBox}
\begin{DefnBox}{Pairwise scoring note}
\texttt{REWRITE\_A}/\texttt{B} presentation order is randomised; signed scores are mapped back to ``A vs.\ B'' after accounting for order.
\end{DefnBox}

\subsection{A.4 Coordinate scorer (LLM backend only)}
When the coordinate backend is LLM-based, sources are rated with the scorer prompt in the persona configuration (integer $S,R\in[-5,5]$, JSON with reasons, then normalise to $[-1,1]$). The runs analysed in the main text use the deterministic HF backend; the LLM configuration branch still documents axis semantics for reviewers.

\section{Appendix B: Automatic metric definitions}
\label{app:metrics}

\begin{DefnBox}{B.1 Novelty N-auto (combined channel)}
Let $c^{\mathrm{d2}}_{\mathrm{rel}}$, $c^{\mathrm{bg}}_{\mathrm{rel}}$, $c^{\mathrm{emb}}_{\mathrm{rel}}$ be genre-relative scores in $[0,1]$. Then
\[
  N_{\mathrm{auto}} = \tfrac{1}{3}\bigl(c^{\mathrm{d2}}_{\mathrm{rel}} + c^{\mathrm{bg}}_{\mathrm{rel}} + c^{\mathrm{emb}}_{\mathrm{rel}}\bigr).
\]
Distinct-2: $z$-score vs.\ genre source mean/std, clip $z\in[-3,3]$, rebound $(z/3+1)/2$. Bigram novelty: offset vs.\ genre baseline with heuristic spread $\sigma{=}0.1$, same clip/rebound. Embedding channel: neutral $0.5$ when SBERT is off; else normalised vs.\ intra-genre mean pairwise distance with clipping.
\end{DefnBox}

\begin{DefnBox}{B.2 Value $V_{\mathrm{auto}}$}
Coherence maps GPT-2 perplexity through a monotone pivot map (default pivot 80). Let $e$ be NLI entailment in $[0,1]$ and $c$ coherence in $[0,1]$. Primary value is the arithmetic mean $(e+c)/2$. Geometric value uses $\sqrt{e\cdot c}$ with a zero guard.
\end{DefnBox}

\begin{DefnBox}{B.3 Creativity variants}
$C_{\mathrm{auto}}=N_{\mathrm{auto}}\cdot V_{\mathrm{arith}}$ (clipped factors); $C_{\mathrm{auto,geom}}=N_{\mathrm{auto}}\cdot V_{\mathrm{geom}}$. Winsorised diagnostics multiply winsorised novelty by arithmetic value in the same analysis protocol as the main tables.
\end{DefnBox}

\section{Appendix C: Statistical robustness}
\label{app:robust}

\subsection{C.1 Pooled quadratic: OLS vs.\ mixed-effects}
Rows match Table~\ref{tab:quadratic}: discrete main, continuous main, and mechanism arms only ($n{=}4620$ with finite $d_H$ and outcomes). Mixed models add a random intercept by source text; optimiser order matches the main analysis script defaults.

\begin{table}[t]
\centering
\begin{tcolorbox}[
  width=\linewidth,
  enhanced, arc=2mm, colback=white, colframe=black,
  title={\texorpdfstring{\textbf{Table C.1.} OLS vs.\ MixedLM on $\beta_{d^2}$ (conflict quadratic)}{\textbf{Table C.1.} OLS vs. MixedLM on beta d2 (conflict quadratic)}},
  fonttitle=\bfseries\small, titlerule=0.55pt, boxrule=0.95pt, boxsep=3mm]
\small
\noindent\resizebox{\linewidth}{!}{%
\begin{tabular}{@{}l l r r r@{}}
\toprule
Metric & Model & $n$ & $\beta_{d^2}$ & $p(d^2)$ \\
\midrule
Creativity auto & MixedLM (L-BFGS, REML) & 4620 & $-0.1111$ & $4.58{\times}10^{-7}$ \\
Creativity auto & OLS (same FE)        & 4620 & $-0.1072$ & $5.97{\times}10^{-5}$ \\
Novelty (combined) & MixedLM (L-BFGS, REML) & 4620 & $0.0908$ & $1.90{\times}10^{-6}$ \\
Novelty (combined) & OLS (same FE)        & 4620 & $0.1143$ & $1.37{\times}10^{-5}$ \\
\bottomrule
\end{tabular}%
}
\end{tcolorbox}
\caption{Quadratic conflict coefficient: sign and significance are stable across OLS and mixed-effects estimators on the pooled main$+$mechanism sample.}
\label{tab:appendix-robust-reg}
\end{table}

\subsection{C.2 T3--T2 gaps under alternative $d_H$ binning}
\begin{table}[t]
\centering
\begin{tcolorbox}[
  width=\linewidth,
  enhanced, arc=2mm, colback=white, colframe=black,
  title={\texorpdfstring{\textbf{Table C.2.} Mean $C_{\mathrm{auto}}$ difference (T3 minus T2) by bucket}{\textbf{Table C.2.} Mean C auto difference (T3 minus T2) by bucket}},
  fonttitle=\bfseries\small, titlerule=0.55pt, boxrule=0.95pt, boxsep=3mm]
\small
\noindent\resizebox{\linewidth}{!}{%
\begin{tabular}{@{}l r r r@{}}
\toprule
Binning scheme & low: T3$-$T2 & mid: T3$-$T2 & high: T3$-$T2 \\
\midrule
Quantile tertiles & $-0.0009$ & $-0.0061$ & $-0.0078$ \\
Equal-width 3 bins & ---       & $-0.0102$ & $-0.0017$ \\
\bottomrule
\end{tabular}%
}
\end{tcolorbox}
\caption{Mean $C_{\mathrm{auto}}$ gaps (T3 minus T2) inside conflict buckets on the same $n{=}4620$ pool; negative values match the main text's ``T3 does not beat T2'' pattern under both binning rules where defined.}
\label{tab:appendix-robust-bin}
\end{table}

\section{Appendix D: Structural collapse risk (exploratory GEE)}
\label{app:collapse}

This appendix reproduces the exploratory \textbf{structural collapse} analysis deferred from the main text's conflict-intensity section.

\subsection{D.1 Indicator and GEE specification}
\begin{DefnBox}{D.1 Collapse event (sample-relative)}
$\mathrm{collapse}{=}1$ iff $\mathrm{NLI}{<}\tau_{\mathrm{NLI}}$ and $\mathrm{coherence}_{\mathrm{auto}}{<}\tau_{\mathrm{coh}}$, with $\tau_{\cdot}$ at the \textbf{same-sample 20th percentiles} on the $n{=}1440$ proprietary T3 discrete rows (same pool as judge-divergence and surprise regressions).
\end{DefnBox}
\noindent
This yields $\tau_{\mathrm{NLI}}{\approx}0.412$, $\tau_{\mathrm{coh}}{\approx}0.628$, and baseline $\hat P(\mathrm{collapse}){\approx}1.9\%$.
Because statsmodels does not expose a standard logit GLMM with Gaussian random intercepts alongside our Gaussian mixed models, we fit a \textbf{binomial GEE} with clustering by source text and exchangeable working correlation (fallback independence); the estimand is marginal (population-averaged) risk.
The genre-inclusive specification $\mathrm{collapse}{\sim}d{+}d^2{+}C(\mathrm{genre})$ shows quasi-separation in narrative/poetry strata; we therefore emphasise the \emph{no-genre} robustness fit for marginal probabilities.
On a fine $d_H$ grid, implied $P(\mathrm{collapse})$ rises from $\approx 3.5\%$ at $d_H{\approx}0.22$ to $\approx 88\%$ at $d_H{\approx}0.88$, with the steepest increase near $d_H{\approx}0.83$; the quadratic term remains significant at $p{\approx}0.003$.
Table~\ref{tab:appendix_collapse_gee} summarises the headline marginal diagnostics.
Full log-odds coefficient tables, the genre-inclusive sensitivity fit, Hessian/convergence metadata, and a median-threshold ($q{=}0.5$) variant are archived with the supplementary regression output for this appendix.

\begin{table}[t]
\centering
\begin{tcolorbox}[
  width=\linewidth,
  enhanced, arc=2mm, colback=white, colframe=black,
  title={\texorpdfstring{\textbf{Table D.1.} Marginal $P(\mathrm{collapse})$ diagnostics (no-$C(\mathrm{genre})$ GEE, $q{=}0.2$)}{\textbf{Table D.1.} Marginal P collapse diagnostics}},
  fonttitle=\bfseries\small, titlerule=0.55pt, boxrule=0.95pt, boxsep=3mm]
\small
\noindent\resizebox{\linewidth}{!}{%
\begin{tabular}{@{}l r@{}}
\toprule
Diagnostic & Value \\
\midrule
Rows ($n$) & 1440 \\
Baseline $\hat P(\mathrm{collapse})$ & ${\approx}1.9\%$ \\
$\hat P(\mathrm{collapse})$ at $d_H{\approx}0.22$ & ${\approx}3.5\%$ \\
$\hat P(\mathrm{collapse})$ at $d_H{\approx}0.88$ & ${\approx}88\%$ \\
$d_H$ at maximum slope of $\hat P$ & ${\approx}0.83$ \\
$p(d_H^2)$ (quadratic term) & ${\approx}0.0032$ \\
\bottomrule
\end{tabular}%
}
\end{tcolorbox}
\caption{Structural-collapse risk: marginal fitted probabilities on $d_H$ (binomial GEE with clustering by source). Log-odds coefficients are omitted because they become large whenever $\hat P$ spans nearly the unit interval; full coefficient tables accompany the supplementary materials.}
\label{tab:appendix_collapse_gee}
\end{table}

\endgroup


\begin{thebibliography}{9}
\bibitem{kahneman}
D.~Kahneman.
\newblock \textit{Thinking, Fast and Slow}.
\newblock Farrar, Straus and Giroux, 2011.

\bibitem{boden}
M.~Boden.
\newblock \textit{The Creative Mind: Myths and Mechanisms}.
\newblock Routledge, 2004.

\bibitem{runco}
M.~Runco and G.~Jaeger.
\newblock The standard definition of creativity.
\newblock \textit{Creativity Research Journal} 24, 1 (2012), 92--96.

\bibitem{cao2025creval}
Q.~Cao, X.~Wang, Y.~Yuan, Y.~Liu, F.~Luo, and R.~Song.
\newblock Evaluating Text Creativity across Diverse Domains: A Dataset and a Large Language Model Evaluator.
\newblock \textit{arXiv preprint arXiv:2505.19236}, 2025.

\bibitem{creval_repo}
Q.~Cao.
\newblock \textit{CrEval}: Public code repository for Evaluating Text Creativity across Diverse Domains.
\newblock GitHub repository, \url{https://github.com/Aman-4-Real/CrEval}. Accessed 2026-05-27.

\bibitem{wang2026crevalimage}
C.~Wang, Z.~Chen, Y.~Wei, T.~Jiang, X.~Wu, F.~Li, W.~Zuo, and H.~Yao.
\newblock {CREval}: An Automated Interpretable Evaluation for Creative Image Manipulation under Complex Instructions.
\newblock \textit{arXiv preprint arXiv:2603.26174}, 2026. \url{https://doi.org/10.48550/arXiv.2603.26174}.

\bibitem{yerkes}
R.~Yerkes and J.~Dodson.
\newblock The relation of strength of stimulus to rapidity of habit-formation.
\newblock \textit{Journal of Comparative Neurology and Psychology} 18 (1908), 459--482.
\end{thebibliography}

\end{document}
